\clearpage 
\section{A simple model of an epidemic}
\subsection{fixed points}
for the fixed points 
\begin{align}
    xky=0 && y\ins{kx-l}=0
\end{align}
so all \(\ins{x,0}\) is a line of fixed points, representing the case where there is no epidemic anymore.
\subsection{trajectories}
\begin{equation}
    \dv{y}{x} = \frac{\dot{y}}{\dot{x}}=\frac{kxy-ly}{-kxy}=-1+\frac{l}{kx}
\end{equation}
integrating 
\begin{equation}
    y=-x+\frac{l}{k}\ln x +C
\end{equation}
so the conserved quantity is 
\begin{equation}
    H = x+y-\frac{l}{k}\ln x
\end{equation}
so we can say 
\begin{equation}
    x+y-\frac{l}{k}\ln x = x_0+y_0-\frac{l}{k}\ln x_0
\end{equation}
therefore 
\begin{equation}
    y = y_0+x_0-x+\frac{l}{k}\ln\ins{\frac{x}{x_0}}
\end{equation}
\subsection{sketch}
\subsubsection{nullclines}
the nullclines for the xdot expression are 
\begin{equation}
    x=0,\quad y=0
\end{equation}
and for ydot
\begin{equation}
    x=\frac{l}{k},\quad y=0
\end{equation}
and for every point \(x,y>0\) we have \(\dot{x}=-kxy<0\), so all the trajectories move left, and it's pretty simple to see that \(\dot{y}>0\) when \(x>l/k\) and vice versa. So the trajectories point left and up to the right of \(x=l/k\) and down to the left of it.
\subsubsection{maximized y}
a maxima happens when \(\dot{y}=0\) so in our case 
\begin{equation}
    x=l/k
\end{equation}
inserting into our expression for \(y\), when \(x_0>l/k\)
\begin{equation}
    y_{max} = y_0+x_0-\frac{l}{k}+\frac{l}{k}\ln\ins{\frac{l/k}{x_0}}
\end{equation}
and when \(x_0<l/k\) \(y\) immediately decreases making \(y_0\) the maximum.
\emph{insert sketch here}
\subsection{Epidemic threshold}
at \(t=0\)
\begin{equation}
    \dot{y}(0)=y_0\ins{kx_0-l}
\end{equation}
and since \(y_0>0\)
\begin{equation}
    \dot{y}(0)>0 \implies kx_0-l>0 \implies x_0 > \frac{l}{k}
\end{equation}
using \(R_0 = kx_0/l\) then the condition is 
\begin{equation}
    R_0 > 1
\end{equation}
lowering the susceptible pool can prevent an outbreak by lowering the ratio \(l/k\), thus lowering \(R_0\) below 1 further.